\documentclass[11pt]{article}
\usepackage[margin=1in]{geometry}
\usepackage{amsmath,amssymb,amsthm}
\usepackage{graphicx}
\usepackage[hidelinks]{hyperref}
\setlength{\emergencystretch}{2em}

\newtheorem{theorem}{Theorem}
\newtheorem{corollary}{Corollary}
\newtheorem{lemma}{Lemma}
\newtheorem{definition}{Definition}

\newcommand{\N}{\mathbb N}
\newcommand{\C}{\mathbb C}
\newcommand{\T}{\mathbb T}
\newcommand{\M}{\mathcal M}
\newcommand{\Lcal}{\mathcal L}
\newcommand{\Kcal}{\mathcal K}

\title{A C*-Annular Diagonal Criterion for Failure of BCP}
\author{Run \texttt{fa\_banach\_001}, lane 19}
\date{June 17, 2026}

\begin{document}
\maketitle

\section*{Source Problem}

The source paper is Jinghao Huang, Karimbergen Kudaybergenov, and Rui Liu,
\emph{The ball-covering property of von Neumann algebras and noncommutative
symmetric spaces}, arXiv:2311.14257.  After proving that a commutative unital
\(C^*\)-algebra \(C(K)\) has BCP/UBCP exactly when \(K\) has a countable
\(\pi\)-basis, the authors ask for a noncommutative version:

\begin{center}
\includegraphics[width=0.92\textwidth]{figures/open_problem_crop.png}
\end{center}

\section*{Proof Intuition}

The quotient-level diagonal criterion says that BCP fails when countably many
centers can be tested on local pieces that nearly see their quotient norm,
while elements of the ideal vanish on those pieces.  In C*-algebras the local
tests usually have the form
\[
        b\longmapsto a_kbc_k
\]
for contractions \(a_k,c_k\) in a multiplier algebra.  The global escaping
element is then built by gluing local pieces over disjoint annuli or corners.

This packet records the C*-version as a theorem.  A second, more concrete
phase-annular corollary covers the common situation where the local pieces
are right annuli \(g_k\), the local escape is a scalar phase \(\lambda_kg_k\),
and one can glue bounded scalar choices into a single element of the ambient
C*-algebra.

\section*{A Localizer Criterion}

Let \(B\) be a \(C^*\)-algebra, let \(I\triangleleft B\) be a closed two-sided
ideal, and let
\[
        q:B\longrightarrow B/I
\]
be the quotient map.  We write \(M(B)\) for the multiplier algebra of \(B\).
If \(B\) is unital, \(M(B)=B\).

\begin{definition}
\label{def:localizer-property}
The quotient \(B/I\) has the C*-localizer diagonal property if for every
sequence \((b_i)_{i\ge1}\subset B\) there are:
\begin{itemize}
\item a map \(\pi:\N\to\N\) with every fiber infinite;
\item contractions \(a_k,c_k\in M(B)\);
\item a null sequence \(\eta_k\downarrow0\);
\item elements \(u_k\in B\);
\item an element \(y\in B\) with \(q(y)\in S_{B/I}\);
\end{itemize}
such that:
\begin{enumerate}
\item \emph{Assigned near-norming}
\[
        \|a_kb_{\pi(k)}c_k\|
        \ge
        \|q(b_{\pi(k)})\|-\eta_k .
\]
\item \emph{Ideal vanishing along fibers}: for every \(d\in I\) and every
\(i\ge1\),
\[
        \lim_{\substack{k\to\infty\\ \pi(k)=i}}\|a_kdc_k\|=0 .
\]
\item \emph{Local gluing}
\[
        \|a_kyc_k-u_k\|\le \eta_k .
\]
\item \emph{Local escape}
\[
        \|u_k-a_kb_{\pi(k)}c_k\|
        \ge
        \|a_kb_{\pi(k)}c_k\|-\eta_k .
\]
\end{enumerate}
\end{definition}

\begin{theorem}
\label{thm:localizer}
If \(B/I\) has the C*-localizer diagonal property, then \(B/I\) fails BCP.
\end{theorem}

\begin{proof}
Let \((\xi_i)_{i\ge1}\subset B/I\) be an arbitrary countable family of
centers.  Choose representatives \(b_i\in B\), so \(q(b_i)=\xi_i\), and
choose the data from Definition~\ref{def:localizer-property}.  Put
\[
        \zeta=q(y)\in S_{B/I}.
\]
Fix \(i\ge1\) and \(d\in I\).  For \(k\) with \(\pi(k)=i\), the map
\[
        L_k:B\to B,\qquad L_k(b)=a_kbc_k
\]
is contractive, and hence
\[
        \|y-b_i-d\|\ge \|a_k(y-b_i-d)c_k\|.
\]
Using local gluing, local escape, and assigned near-norming gives
\begin{align*}
        \|a_k(y-b_i-d)c_k\|
        &\ge
        \|u_k-a_kb_ic_k\|-\|a_kyc_k-u_k\|-\|a_kdc_k\|\\
        &\ge
        \|a_kb_ic_k\|-2\eta_k-\|a_kdc_k\|\\
        &\ge
        \|\xi_i\|-3\eta_k-\|a_kdc_k\|.
\end{align*}
Letting \(k\to\infty\) through the fiber \(\pi^{-1}(\{i\})\), and using ideal
vanishing, gives
\[
        \|y-b_i-d\|\ge \|\xi_i\|.
\]
Taking the infimum over \(d\in I\), we obtain
\[
        \|\zeta-\xi_i\|\ge \|\xi_i\|.
\]
Thus the unit vector \(\zeta\) is outside every open ball
\[
        B(\xi_i,\|\xi_i\|).
\]
Since the countable family of centers was arbitrary, \(B/I\) fails BCP.
\end{proof}

\section*{A Phase-Annular Corollary}

The next corollary is often the easiest form to check.  It packages the
Hahn--Banach phase trick and the auxiliary unit-norm check into the
hypotheses.

\begin{theorem}[phase-annular criterion]
\label{thm:phase-annular}
Let \(B\) be a unital \(C^*\)-algebra and let \(I\triangleleft B\) be a
proper closed two-sided ideal.  Suppose that for every sequence
\((b_i)_{i\ge1}\subset B\), after setting \(b_0=1_B\), there are:
\begin{itemize}
\item a map \(\pi:\N\to\{0,1,2,\ldots\}\) with every fiber infinite;
\item contractions \(g_k\in B\);
\item a null sequence \(\eta_k\downarrow0\);
\item for every scalar sequence \(\lambda=(\lambda_k)\in\ell_\infty\) with
\(\|\lambda\|_\infty\le1\), an element \(\Phi(\lambda)\in B\);
\end{itemize}
such that:
\begin{enumerate}
\item \emph{Assigned right-annulus near-norming}
\[
        \|b_{\pi(k)}g_k\|
        \ge
        \|q(b_{\pi(k)})\|-\eta_k,\qquad k\ge1.
\]
\item \emph{Ideal vanishing}: for every \(d\in I\) and every
\(i\ge0\),
\[
        \lim_{\substack{k\to\infty\\ \pi(k)=i}}\|dg_k\|=0 .
\]
\item \emph{Scalar gluing}
\[
        \|q(\Phi(\lambda))\|\le1,
        \qquad
        \|\Phi(\lambda)g_k-\lambda_kg_k\|\le\eta_k
        \quad(k\ge1).
\]
\end{enumerate}
Then \(B/I\) fails BCP.
\end{theorem}

\begin{proof}
Let \((\xi_i)_{i\ge1}\subset B/I\) be an arbitrary countable family, and
choose representatives \(b_i\in B\).  Add \(b_0=1_B\) and choose the data in
the statement.

For each \(k\) with \(\pi(k)\ge1\), choose \(\lambda_k\in\T\) such that
\[
        \|\lambda_kg_k-b_{\pi(k)}g_k\|\ge \|b_{\pi(k)}g_k\|.
\]
Indeed, if \(b_{\pi(k)}g_k\ne0\), take a norm-one functional \(\varphi\) with
\(\varphi(b_{\pi(k)}g_k)=\|b_{\pi(k)}g_k\|\), and choose \(\lambda_k\in\T\)
so that \(\operatorname{Re}(\lambda_k\varphi(g_k))\le0\).  If
\(b_{\pi(k)}g_k=0\), any \(\lambda_k\in\T\) works.  For \(k\) with
\(\pi(k)=0\), put \(\lambda_k=1\).

Let
\[
        y=\Phi((\lambda_k)_{k\ge1}).
\]
The scalar gluing hypothesis gives \(\|q(y)\|\le1\).  We claim that
\(\|q(y)\|=1\).  Fix \(d\in I\) and let \(k\to\infty\) through the fiber
\(\pi^{-1}(\{0\})\).  Since \(b_0=1_B\),
\[
        \|g_k\|=\|b_0g_k\|\ge \|q(1_B)\|-\eta_k=1-\eta_k.
\]
Therefore
\[
        \|y-d\|
        \ge
        \|(y-d)g_k\|
        \ge
        \|\lambda_kg_k\|-\|yg_k-\lambda_kg_k\|-\|dg_k\|
        \ge
        1-2\eta_k-\|dg_k\|.
\]
Letting \(k\to\infty\) along that fiber gives \(\|y-d\|\ge1\), and then
taking the infimum over \(d\in I\) gives \(\|q(y)\|\ge1\).  Hence
\[
        q(y)\in S_{B/I}.
\]

Now fix \(i\ge1\) and \(d\in I\).  Let \(k\to\infty\) through
\(\pi^{-1}(\{i\})\).  Then
\begin{align*}
        \|y-b_i-d\|
        &\ge \|(y-b_i-d)g_k\|\\
        &\ge \|\lambda_kg_k-b_ig_k\|-\|yg_k-\lambda_kg_k\|-\|dg_k\|\\
        &\ge \|b_ig_k\|-\eta_k-\|dg_k\|\\
        &\ge \|q(b_i)\|-2\eta_k-\|dg_k\|.
\end{align*}
Letting \(k\to\infty\) along the fiber and then taking the infimum over
\(d\in I\), we obtain
\[
        \|q(y)-\xi_i\|\ge \|\xi_i\|.
\]
Thus \(q(y)\) is a unit vector outside every proposed ball
\(B(\xi_i,\|\xi_i\|)\), and \(B/I\) fails BCP.
\end{proof}

\section*{Recovering Existing Diagonal Packets}

\subsection*{Sigma-Unital Multiplier Coronas}

Let \(A\) be a nonunital \(\sigma\)-unital \(C^*\)-algebra.  In the
sigma-unital corona packet, one takes
\[
        B=M(A),\qquad I=A.
\]
An increasing positive contractive approximate unit \((e_n)\) with
\(e_me_n=e_n\) for \(m\ge n\) supplies inner annuli
\[
        g_k=e_{m_k}-e_{n_k}
\]
and larger outer annuli
\[
        h_k=e_{b_k}-e_{a_k},\qquad h_kg_k=g_k,\qquad h_jg_k=0\ (j\ne k).
\]
Finite annulus recovery gives
\[
        \|T_{\pi(k)}g_k\|\approx \|q(T_{\pi(k)})\|,
\]
and \(A\)-elements vanish on far-out annuli.  For
\(\lambda\in\ell_\infty\), strict convergence gives the scalar gluing element
\[
        \Phi(\lambda)=\sum_k\lambda_kh_k\in M(A),
\]
with \(\|\Phi(\lambda)\|\le\|\lambda\|_\infty\) and
\(\Phi(\lambda)g_k=\lambda_kg_k\).  Thus Theorem~\ref{thm:phase-annular}
is exactly the C*-annular checklist behind the sigma-unital corona theorem.

\subsection*{Hilbert-Module Coronas}

Let \(E\) be a Hilbert \(C^*\)-module with a BCP block system.  In the
Hilbert-module block packet, one takes
\[
        B=\Lcal(E),\qquad I=\Kcal(E).
\]
The localizers in Theorem~\ref{thm:localizer} are
\[
        a_k=1,\qquad c_k=e_{F_k},
\]
where \(e_{F_k}\) is a fresh compact block projection.  Tail norm recovery
gives assigned near-norming, compact-vanishing gives ideal vanishing, fresh
range escape supplies the local pieces \(u_k=v_k\), and strict gluing gives a
single operator \(y=V\) with \(Ve_{F_k}=v_k\).  Hence the Hilbert-module block
criterion is an instance of the localizer theorem.

\subsection*{Ordinary Product Quotients}

For an ordinary \(c_0\)-product quotient
\[
        \prod_n A_n\Big/\bigoplus\nolimits_{c_0}A_n
\]
with unital \(A_n\)'s, disjoint infinite coordinate sets \(T_k\) give central
projections \(g_k=1_{T_k}\).  Elements of the \(c_0\)-sum vanish on these
sets as \(k\to\infty\), after the sets are chosen pairwise disjoint.  Local
escape elements can be glued coordinatewise.  This is the coordinate version
of the same annular principle.  For more general ideals \(\mathcal I\), the
full reduced-product packet is better viewed through the broader
quotient-level asymptotic criterion, because local restrictions naturally
land in local reduced quotients rather than in a single ambient C*-algebra.

\section*{Verification and Search Bounds}

The run indexes were searched for arXiv:2311.14257 together with annular
diagonal, C*-localizer, multiplier corona, Hilbert-module block criterion,
asymptotic diagonal, and quotient diagonal keywords.  Existing packets
contained the Banach quotient criterion and concrete C*-diagonal proofs, but
no standalone C*-localizer/phase-annular criterion.

The proof of Theorem~\ref{thm:localizer} is a direct quotient-norm diagonal
estimate.  The phase-annular theorem is a self-contained specialization using
the Hahn--Banach phase choice and an auxiliary \(b_0=1\) fiber to force the
glued quotient element to have norm one.

\section*{Human Review Focus}

Check the localizer hypotheses in Definition~\ref{def:localizer-property} and
the scalar gluing hypothesis in Theorem~\ref{thm:phase-annular}.  These are
the reusable C*-checklists; once they are verified in an application, the BCP
failure proof is formal.

\begin{thebibliography}{9}

\bibitem{HKL}
Jinghao Huang, Karimbergen Kudaybergenov, and Rui Liu,
\emph{The ball-covering property of von Neumann algebras and noncommutative
symmetric spaces},
arXiv:2311.14257, 2023.

\end{thebibliography}

\end{document}
